\documentclass[12pt]{memoir}

% QIC 890 course project --- LaTeX template.

%-----------------------------------------------------------------------------%
% Packages

\usepackage[T1]{fontenc}
\usepackage{mathpazo}
\usepackage{amsmath, amssymb, amsthm, mathtools}
\usepackage{graphicx}
\usepackage{needspace}
\usepackage{tikz}
\usepackage[colorlinks=true, linkcolor=blue, citecolor=blue,
  urlcolor=blue]{hyperref}

\frenchspacing
\sloppybottom

%-----------------------------------------------------------------------------%
% Page layout (the same as the course notes)

\settypeblocksize{8.5in}{5.5in}{*}
\setlrmarginsandblock{1.25in}{1.25in}{*}
\setulmargins{1.25in}{*}{*}
\checkandfixthelayout
\linespread{1.15}

\pagestyle{plain}
\setsecnumdepth{subsection}

% There are no chapters, so sections are numbered 1, 2, 3, ...
% The bibliography is a section rather than a chapter.
\renewcommand{\thesection}{\arabic{section}}
\renewcommand{\bibsection}{\section*{References}}

%-----------------------------------------------------------------------------%
% Theorem-like environments, numbered within sections

\theoremstyle{definition}
\newtheorem{definition}{Definition}[section]
\newtheorem{example}[definition]{Example}

\theoremstyle{plain}
\newtheorem{theorem}[definition]{Theorem}
\newtheorem{lemma}[definition]{Lemma}
\newtheorem{corollary}[definition]{Corollary}
\newtheorem{proposition}[definition]{Proposition}

\theoremstyle{remark}
\newtheorem{remark}[definition]{Remark}

\numberwithin{equation}{section}

%-----------------------------------------------------------------------------%
% Macros from the course notes

% Dirac notation, norms, and absolute values
\DeclarePairedDelimiter\bra{\langle}{\rvert}
\DeclarePairedDelimiter\ket{\lvert}{\rangle}
\DeclarePairedDelimiterX\braket[2]{\langle}{\rangle}{#1 \delimsize\vert #2}
\DeclarePairedDelimiter\norm{\lVert}{\rVert}
\DeclarePairedDelimiter\abs{\lvert}{\rvert}

\newcommand{\yes}{\mathrm{yes}}
\newcommand{\no}{\mathrm{no}}

% Acceptance and rejection probabilities of a circuit. The trailing optional
% argument is the input the circuit is run on, and is left off for circuits
% that take no input:
%   \accepts{Q_x}     -->  Pr(Q_x accepts)
%   \accepts{Q_n}[x]  -->  Pr(Q_n accepts x)
\NewDocumentCommand{\accepts}{m o}{\Pr(#1 \;\text{accepts}\IfValueT{#2}{\;#2})}
\NewDocumentCommand{\rejects}{m o}{\Pr(#1 \;\text{rejects}\IfValueT{#2}{\;#2})}

% Complexity classes (\L and \P are taken by LaTeX, hence \Lclass and \Pclass).
\newcommand{\Lclass}{\mathrm{L}}
\newcommand{\Pclass}{\mathrm{P}}
\newcommand{\NP}{\mathrm{NP}}
\newcommand{\coNP}{\mathrm{coNP}}
\newcommand{\BPP}{\mathrm{BPP}}
\newcommand{\PP}{\mathrm{PP}}
\newcommand{\PSPACE}{\mathrm{PSPACE}}
\newcommand{\EXP}{\mathrm{EXP}}
\newcommand{\NEXP}{\mathrm{NEXP}}
\newcommand{\MA}{\mathrm{MA}}
\newcommand{\AM}{\mathrm{AM}}
\newcommand{\BQP}{\mathrm{BQP}}
\newcommand{\PostBQP}{\mathrm{PostBQP}}
\newcommand{\QMA}{\mathrm{QMA}}
\newcommand{\QCMA}{\mathrm{QCMA}}
\newcommand{\QIP}{\mathrm{QIP}}
\newcommand{\SharpP}{\mathrm{\#P}}
\newcommand{\GapP}{\mathrm{GapP}}

% Gates and quantum notation
\newcommand{\reg}[1]{\mathsf{#1}}
\newcommand{\CNOT}{\mathrm{CNOT}}
\newcommand{\TOF}{\mathrm{TOF}}
\newcommand{\Id}{\mathbb{1}}
\newcommand{\tr}{\operatorname{Tr}}
\newcommand{\dnorm}[1]{\left\lVert #1 \right\rVert_{\diamond}}

% Karp reductions
\newcommand{\redPT}{\leq^{\mathrm{p}}_{\mathrm{m}}}

% Promise problem display:
%   \begin{promiseproblem}{Name of the problem (ABBREVIATION)}
%     \item[Input] ...
%     \item[Yes]   ...
%     \item[No]    ...
%   \end{promiseproblem}
\newenvironment{promiseproblem}[1]
  {\par\medskip%
   \needspace{6\baselineskip}%
   \noindent\textsc{#1}\par\nopagebreak[4]%
   \begin{list}{}{%
     \setlength{\leftmargin}{5em}%
     \setlength{\labelwidth}{4em}%
     \setlength{\labelsep}{1em}%
     \setlength{\itemindent}{0pt}%
     \setlength{\topsep}{6pt}%
     \setlength{\partopsep}{0pt}%
     \setlength{\parsep}{0pt}%
     \setlength{\itemsep}{2pt}%
     \renewcommand{\makelabel}[1]{\hfil\emph{##1}}%
   }}
  {\end{list}\medskip}

%-----------------------------------------------------------------------------%
% Title

\newcommand{\reporttitle}{Title of the Report}
\newcommand{\reportauthor}{Your Name}

\begin{document}

\begin{center}
  {\small QIC 890: Quantum Complexity Theory --- Course Project, Fall 2026}
  \\[8mm]
  {\fontsize{24}{28}\selectfont \reporttitle}\\[6mm]
  {\large \reportauthor}
\end{center}

\vspace{6mm}

%-----------------------------------------------------------------------------%

Opening paragraphs.

%-----------------------------------------------------------------------------%
\section{Background}

Background material and definitions. Results proved in the course notes can be
used without proof, referred to by number, for instance Theorem~6.17.

\begin{definition}
  \label{def:example}
  A sample definition, in which a new \emph{term} is introduced.
\end{definition}

%-----------------------------------------------------------------------------%
\section{Main result}

\begin{theorem}[An important theorem]
  \label{thm:main}
  A sample theorem statement including an optional theorem name, a reference
  to Definition~\ref{def:example}, and a citation~\cite{Kitaev--Shen--Vyalyi}.
\end{theorem}

\begin{proof}
  A sample proof, with a numbered display:
  \begin{equation}
    \label{eq:example}
    \bra{\psi} \Pi \ket{\psi} = \tr\bigl(\Pi \ket{\psi}\bra{\psi}\bigr).
  \end{equation}
  Equation~\eqref{eq:example} is referred to by its number.
\end{proof}

%-----------------------------------------------------------------------------%
\section{Further remarks}

Extensions, open problems, and pointers to related work.

%-----------------------------------------------------------------------------%
% References. Anything other than the course notes must be cited. Use
% BibLaTeX or BibTeX to generate automatically.

\begin{thebibliography}{9}

\bibitem{Kitaev--Shen--Vyalyi}
  A.~Yu. Kitaev, A.~H. Shen, and M.~N. Vyalyi.
  \emph{Classical and Quantum Computation}.
  Graduate Studies in Mathematics, volume~47.
  American Mathematical Society, 2002.

\end{thebibliography}

\end{document}
